\chapter{اللقاء السابع عشر: تابع تفسير سورة البقرة (٦)}

\section{مقدمة ومناقشة واجب (أم)}
السلام عليكم ورحمة الله وبركاته. الحمد لله رب العالمين، وأشهد أن لا إله إلا الله وحده لا شريك له، وأشهد أن محمداً عبده ورسوله. اللهم صل على محمد وعلى آل محمد كما صليت على آل إبراهيم إنك حميد مجيد، اللهم بارك على محمد وعلى آل محمد كما باركت على آل إبراهيم إنك حميد مجيد.

قال الله تبارك وتعالى: \begin{ayah}يَا أَيُّهَا النَّاسُ قَدْ جَاءَتْكُم مَّوْعِظَةٌ مِّن رَّبِّكُمْ وَشِفَاءٌ لِّمَا فِي الصُّدُورِ وَهُدًى وَرَحْمَةٌ لِّلْمُؤْمِنِينَ * قُلْ بِفَضْلِ اللَّهِ وَبِرَحْمَتِهِ فَبِذَٰلِكَ فَلْيَفْرَحُوا هُوَ خَيْرٌ مِّمَّا يَجْمَعُونَ\end{ayah}. والحمد لله على نعمة الإسلام، والحمد لله على نعمة الهداية بالقرآن، والحمد لله الذي أنزل على عبده الكتاب ولم يجعل له عوجاً، والحمد لله أن جعلنا ممن يقبلون الهدى الذي جاء به النبي محمد ﷺ وهو القرآن. نسأل الله أن يجعلنا ممن فَقُه في دين الله ونفعه الله به فَعَلِم وعَلَّم.

صباح الخير يا شباب. أهلاً وسهلاً ومرحباً بطلبة العلم الكرام. أسأل الله أن يبارك فيكم جميعاً وأن يحفظكم وأن يهديكم ويسددكم، وأن يثبت قلوبكم على حب العلم وأن يهديكم إلى العلم النافع الصالح. وكثيراً ما أحب من طالب العلم أن يستشعر أنه بخير، أنه يحب الإسلام يحب الله يحب رسوله، يحب العلم يحب القرآن ويدخر وقته لذلك. هذا والله خير ما يمكن أن يشتغل به إنسان؛ \textit{رسول من الله يتلو صحفاً مطهرة}، هذه الثنائية كل من قَبِلها وعمل بها واهتدى بها فهو من خير البرية.

نبدأ بمناقشة الواجبات الخاصة بالدرس الماضي، وكان منها تلخيص أحوال حرف (أم).

\subsection{أوجه (أم) في لغة العرب}
لخّص أحد الطلاب من كتاب «مغني اللبيب» لـ \rawi{ابن هشام الأنصاري} أن لـ (أم) أوجهاً:
\begin{itemize}
    \item \textbf{أم المتصلة (المعادلة):} وهي التي يقع ما بعدها وما قبلها لا يُستغنى بأحدهما عن الآخر، وتُسبق إما بهمزة تسوية (مثل: \textit{سواء عليهم أأنذرتهم أم لم تنذرهم})، أو بهمزة يطلب بها التعيين (مثل: \textit{أزيد في الدار أم عمرو}).
    \item \textbf{أم المنقطعة:} وهي التي بمعنى (بل) والهمزة المقدرة للاستفهام (غالباً الإنكاري)، وتقع بعد خبر محض (مثل: \textit{تنزيل الكتاب لا ريب فيه من رب العالمين * أم يقولون افتراه})، أو بعد همزة ليست للاستفهام (مثل: \textit{ألهم أرجل يمشون بها أم لهم أيد يبطشون بها})، أو بعد استفهام بغير الهمزة (مثل: \textit{هل يستوي الأعمى والبصير أم هل تستوي الظلمات والنور}).
    \item \textbf{أم بمعنى أل (للتعريف):} وهي لغة قبائل طيئ وحمير في اليمن (مثل: \textit{ليس من امبر امصيام في امسفر}، أي ليس من البر الصيام في السفر).
    \item \textbf{أم الزائدة:} ذكر \rawi{أبو زيد} وغيره أنها قد تأتي زائدة (مثل توجيههم لقوله: \textit{أفلا تبصرون أم أنا خير} أي أفلا تبصرون أنا خير).
\end{itemize}

\begin{fawaid}[بطلان الزيادة في القرآن بلا معنى]
مذهب المحققين كالإمام الطبري وشيخ الإسلام ابن تيمية عدم الموافقة على أن يكون في كتاب الله حرف "زائد" لا معنى له (بحيث يكون وجوده وعدمه سواء). وما أُطلق عليه "زائد" من قبل بعض اللغويين مقصودهم أنه يمكن استقامة الجملة إعرابياً دونه، ولكن وجوده يضيف معنى وتأكيداً لا يتم بغيره. فدعوى الزيادة المجردة تُنكر.
\end{fawaid}

\separator

\section{تأويل (ود كثير من أهل الكتاب لو يردونكم من بعد إيمانكم كفارا)}
\isnad{قال أبو جعفر:} «وقد صرح هذا القول من قول الله جل ثناؤه... أن خطابه بجميع هذه الآيات من قوله \textit{يا أيها الذين آمنوا لا تقولوا راعنا}... إنما هو خطاب منه للمؤمنين من أصحابه، وعتاب منه لهم، ونهي عن انتصاح اليهود ونظرائهم من أهل الشرك وقبول آرائهم في شيء من أمور دينهم... فإن اليهود والمشركين ما يودون أن ينزل عليكم من خير من ربكم، ولكن كثيراً منهم ودوا أنهم يردونكم من بعد إيمانكم كفاراً حسداً من عند أنفسهم... من بعد ما تبين لهم الحق في أمر محمد ﷺ وأنه نبي إليهم وإلى خلقه كافة».

ورَدَّ \rawi{الطبري} على من زعم أن قوله (ود كثير من أهل الكتاب) يعني به رجلاً واحداً وهو كعب بن الأشرف (بحجة كثرة ماله أو عزه)، مبيناً أنه لا دلالة في الآية على أن المراد به واحد دون جماعة كثيرة.
وأصّل قاعدة: \textbf{«لا يجوز صرف تأويل الآية وإحالة دليل ظاهرها إلى غير الغالب في الاستعمال»}.

\subsection{تأويل (حسدا من عند أنفسهم)}
\isnad{قال أبو جعفر:} «والحسد إذًا منصوب على غير النعت للكفار، ولكن على وجه المصدر الذي يأتي خارجاً من معنى الكلام... لأن في قوله \textit{ود كثير من أهل الكتاب لو يردونكم...} معنى حسدكم أهل الكتاب على ما أعطاكم الله التوفيق وهب لكم من الرشاد لدينه... وخصكم به من أن جعل رسوله إليكم رجلاً منكم... فكان قوله (حسداً) مصدراً من ذلك المعنى».

\isnad{وأما قوله (من عند أنفسهم):} «فإنه يعني بذلك من قِبَل أنفسهم... إعلاماً منه لهم أنهم لم يؤمروا بذلك في كتابهم... وأنهم يأتون ما يأتون من ذلك على علم منهم بنهي الله إياهم عنه».

\begin{fawaid}[الحسد من عند أنفسهم وتبرئة الوحي]
في قوله تعالى: \textit{حَسَدًا مِّنْ عِندِ أَنفُسِهِم} دلالتان عظيمتان: 
\begin{itemize}
    \item \textbf{الأولى:} أن أبواب الباطل لا تقتصر على الجهل، بل قد يكون دافعها مرض القلب كالحسد والكبر (كما حسد اليهود النبي ﷺ لأن النبوة انتقلت للعرب فلم يرضوا أن يكونوا أتباعاً). 
    \item \textbf{الثانية:} تبرئة التوراة من أن تكون قد أمرتهم بذلك، فهو حسد نابع من عند أنفسهم ولم يأمرهم به كتابهم، فهذا من فعلهم هم وليس من تشريع الله لهم. وهذا أصل في التفريق بين الدين الإسلامي الصافي وبين أفعال بعض المنتسبين إليه الذين قد يخالفونه بأفعالهم. ومن أجلّ ما يقوم به طالب العلم الذب عن الشريعة بنفي أفعال المنتسبين إليها التي لا تمثلها.
\end{itemize}
\end{fawaid}

\separator

\section{تأويل (من بعد ما تبين لهم الحق)}
\isnad{قال أبو جعفر:} «يعني جل ثناؤه بقوله \textit{من بعد ما تبين لهم الحق} أي من بعد ما تبين لهؤلاء الكثير من أهل الكتاب... الحق في أمر محمد ﷺ وما جاء به من عند ربه والملة التي دعا إليها، فأضاء لهم أن ذلك الحق الذي لا يمترون فيه».

\begin{fawaid}[عقوبة رد الحق بعد تبينه]
تُبين الآية أن كفر هؤلاء لم يكن عن جهل بل عناد (من بعد ما تبين لهم الحق). وكل من تبين له الحق فعانده أو أعرض عنه كبراً أو حسداً، فمن عقوبته المغلظة أن يُعْمَى عن الحق ويُصرف عنه لاحقاً، كما قال تعالى: \textit{سَأَصْرِفُ عَنْ آيَاتِيَ الَّذِينَ يَتَكَبَّرُونَ فِي الْأَرْضِ بِغَيْرِ الْحَقِّ وَإِن يَرَوْا كُلَّ آيَةٍ لَّا يُؤْمِنُوا بِهَا}، وقال: \textit{وَنُقَلِّبُ أَفْئِدَتَهُمْ وَأَبْصَارَهُمْ كَمَا لَمْ يُؤْمِنُوا بِهِ أَوَّلَ مَرَّةٍ}.
\end{fawaid}

\separator

\section{تأويل (فاعفوا واصفحوا حتى يأتي الله بأمره)}
\isnad{قال أبو جعفر:} «يعني جل ثناؤه بقوله \textit{فاعفوا} فتجاوزوا عما كان منهم من إساءة... \textit{واصفحوا} عما كان منهم من جهل في ذلك \textit{حتى يأتي الله بأمره} فيحدث لكم من أمره ما يشاء ويقضي فيهم بما يريد... فقضى فيهم بعد ذلك... \textit{قاتلوا الذين لا يؤمنون بالله ولا باليوم الآخر}... فنسخ العفو جل ثناؤه عنهم والصفح بفرض قتالهم».

\begin{fawaid}[العفو والصفح بين النسخ والإحكام]
ذكر الطبري وأكثر المفسرين المتقدمين أن هذه الآية (فاعفوا واصفحوا) نُسخت بآيات القتال (السيف). وبعض أهل العلم يرى أن هذه الآية محكمة، وأن الحكم يختلف بحسب الظرف؛ ففي حال الاستضعاف يُستدعى العفو والصفح، وفي حال القوة والمناجزة تُستدعى آيات القتال. والنسخ عند المتقدمين أوسع من إزالة الحكم بالكلية، بل يشمل تخصيص الحال.
\end{fawaid}

\separator

\section{تأويل (إن الله على كل شيء قدير)}
\isnad{قال أبو جعفر:} «قد دللنا فيما مضى على معنى القدير وأنه القوي. فمعنى الآية هاهنا: إن الله على كل ما يشاء ويريد بالذين وصفت لكم أمرهم من أهل الكتاب وغيرهم قدير... لا يتعذر عليه شيء أراده، ولا يمتنع عليه أمر شاء قضاءه، لأن له الخلق والأمر».

\begin{fawaid}[الرد على المعتزلة في قدرة الله]
قوله تعالى: \textit{إِنَّ اللَّهَ عَلَىٰ كُلِّ شَيْءٍ قَدِيرٌ} عموم لا يدخله التخصيص. وفيه رد على المعتزلة الذين زعموا أن الله لا يقدر على أفعال العباد (لأنه لو كان قادراً عليها ثم عذبهم لكان ظلماً بزعمهم!)، فجعلوا الآية عاماً مُخَصَّصاً، وهذا ضلال. فالله قادر على كل شيء، سواء شاءه فخلقه، أم لم يشأه، كما قال: \textit{بَلَىٰ قَادِرِينَ عَلَىٰ أَن نُّسَوِّيَ بَنَانَهُ}. 
ومن الخطأ الذي يقع فيه بعض الخطباء حصر العبارة في قولهم (إن الله على ما يشاء قدير)، وهو وإن كان المعنى صحيحاً وموجوداً في بعض النصوص، إلا أنه قاصر كإطلاق عام عن إثبات قدرته المطلقة على ما شاء وما لم يشأ. فالمحكم هو: (إن الله على كل شيء قدير).
\end{fawaid}

\separator

\section{تأويل (وأقيموا الصلاة وآتوا الزكاة وما تقدموا لأنفسكم من خير تجدوه عند الله)}
\isnad{قال أبو جعفر:} «وإنما أمرهم جل ثناؤه في هذا الموضع بما أمرهم به من إقام الصلاة وإيتاء الزكاة وتقديم الخيرات لأنفسهم، ليتطهروا بذلك من الخطأ الذي سلف منهم في استنصاحهم اليهود وركون ما كان ركن منهم إليهم... وفي تقديم الخيرات إدراك الفوز برضوان الله».

\begin{fawaid}[الثبات على العمل الصالح وقت الفتن]
بعد أن بيّن الله مكر اليهود والمشركين، أمر المؤمنين بقوله: \textit{وَأَقِيمُوا الصَّلَاةَ وَآتُوا الزَّكَاةَ}، ليرشدهم إلى أن أعظم ما يُواجه به كيد الأعداء ومكرهم هو الثبات على الإيمان والعمل الصالح والازدياد منه. فالمؤمن لا يقعده الحزن على كيد الكائدين عن طاعة الله، بل يستمسك بدينه.
\end{fawaid}

\isnad{وأما قوله:} \begin{ayah}إِنَّ اللَّهَ بِمَا تَعْمَلُونَ بَصِيرٌ\end{ayah} \isnad{قال الطبري:} «هذا الكلام وإن كان خرج مخرج الخبر فإن فيه وعداً ووعيداً وأمراً وزجراً. وذلك أنه أعلم القوم أنه بصير بجميع أعمالهم... ليجدوا في طاعته... وليحذروا معصيته».

\separator

\section{تأويل (وقالوا لن يدخل الجنة إلا من كان هودا أو نصارى تلك أمانيهم)}
\isnad{قال أبو جعفر:} «يعني جل ثناؤه... وقالت اليهود والنصارى: لن يدخل الجنة. فإن قال قائل: وكيف جمع اليهود والنصارى في هذا الخبر مع اختلاف مقالة الفريقين؟... قيل: إنما عنى به: وقالت اليهود لن يدخل الجنة إلا من كان هوداً، وقالت النصارى لن يدخل الجنة إلا من كان نصارى، ولكن جمع الفريقان في الخبر عنهما فقيل: وقالوا...».

\isnad{وأما قوله:} \begin{ayah}تِلْكَ أَمَانِيُّهُمْ\end{ayah} \isnad{قال الطبري:} «أماني منهم يتمنونها على الله بغير حق ولا حجة ولا برهان، ولا يقين علم بصحة ما يدعون، ولكن بادعاء الأباطيل وأماني النفوس الكاذبة».

\separator

\section{تأويل (قل هاتوا برهانكم إن كنتم صادقين)}
\isnad{قال أبو جعفر:} «يقول لمحمد ﷺ: قل للزاعمين أن الجنة لا يدخلها إلا من كان هوداً أو نصارى دون غيرهم من سائر البشر: هاتوا حجتكم على ما تزعمون من ذلك فنسلم لكم دعواكم... والبرهان هو البيان والحجة والبينة».
وبيّن \rawi{الطبري} أن هذا أمر تعجيزي وتكذيبي، لأنهم لا يملكون برهاناً، كما في قوله: \textit{فأتوا بسورة من مثله}.

\separator

\section{تأويل (بلى من أسلم وجهه لله وهو محسن)}
\isnad{قال أبو جعفر:} «يعني بقوله جل ثناؤه: بلى من أسلم... ليس الأمر كما قال الزاعمون لن يدخل الجنة إلا من كان هوداً أو نصارى، ولكن من أسلم وجهه لله وهو محسن فهو الذي يدخلها وينعم فيها».

\isnad{ومعنى إسلام الوجه:} «التذلل لطاعته والإذعان لأمره. وأصل الإسلام الاستسلام... وإنما سمي المسلم مسلماً لخضوع جوارحه بطاعة ربه... وخص الله جل ثناؤه بالخبر عمن أخبر عنه... بإسلام وجهه له دون سائر جوارحه، لأن أكرم أعضاء بني آدم وجوارحه وجهه... فإذا خضع لشيء وجهه الذي هو أكرم جزء بجسده عليه، فغيره من أجزاء جسده أحرى أن يكون قد خضع له... فاكتفى بذكر الوجه من ذكر جسده».

\isnad{ومعنى محسن:} «يقول وهو محسن في إسلامه لله جسده... الا من أخلص طاعته وعبادته لله».

\separator

\section{تأويل (وقالت اليهود ليست النصارى على شيء وقالت النصارى ليست اليهود على شيء وهم يتلون الكتاب)}
ذكر \rawi{الطبري} الأخبار عن \rawi{ابن عباس} وغيره أنها نزلت حين قدم نصارى نجران فتنازعوا مع أحبار اليهود عند رسول الله ﷺ وكفّر كل فريق منهم الآخر.
\isnad{قال الطبري:} «فاخبر جل ثناؤه أن كل فريق منهم قال ما قال من ذلك على علم منهم أنهم فيما قالوه مبطلون... وتلاوة كل واحد من الفريقين كتابه الذي يشهد على كذبه في قيله ذلك». فالتوراة تصدق عيسى والإنجيل يصدق موسى.

\isnad{وأما قوله:} \begin{ayah}كَذَٰلِكَ قَالَ الَّذِينَ لَا يَعْلَمُونَ مِثْلَ قَوْلِهِمْ\end{ayah}
أبقى \rawi{الطبري} لفظ (الذين لا يعلمون) على عمومه، سواء كانوا مشركي العرب أو أمماً سابقة بلا كتاب، لأنهم شاركوهم في الجهل والتكذيب.
وبيّن غرض الآية: «وإنما قصد الله جل ثناؤه... إلى إعلام المؤمنين أن اليهود والنصارى قد أتوا من قيل الباطل... وهم أهل كتاب يعلمون أنهم فيما يقولون مبطلون... مثل الذي قاله أهل الجهل بالله وكتبه ورسله... وهذه الآية تنبئ عن أن من أتى شيئاً من معاصي الله على علم منه... فمصيبته في دينه أعظم من مصيبة من أتى ذلك جاهلاً به».

\begin{fawaid}[أمانة العلم وتلاوة الكتاب]
وبّخ الله اليهود والنصارى على اختلافهم وتكفير بعضهم لبعض مع أنهم يقرأون الكتاب المصدق للأنبياء (\textit{وَهُمْ يَتْلُونَ الْكِتَابَ})، وسوّى بينهم وبين المشركين الجاهلين الذين لا كتاب لهم (\textit{كَذَٰلِكَ قَالَ الَّذِينَ لَا يَعْلَمُونَ مِثْلَ قَوْلِهِمْ}). وهذا يوجب على طالب العلم أن يكون "وقّافاً" عند كتاب الله عاملاً به، وإلا كان هو والجاهل سواء.
\end{fawaid}

\separator

\section{تأويل (ومن أظلم ممن منع مساجد الله أن يذكر فيها اسمه وسعى في خرابها)}
\isnad{قال أبو جعفر:} «أي: وأي امرئ أشد تعدياً وجرأة على الله وخلافاً لأمره من امرئ منع مساجد الله أن يعبد الله فيها...».

واختلف المفسرون في المَعْنِيّ بالآية: فقيل هم النصارى (سعوا مع بختنصر في خراب بيت المقدس)، وقيل هم بختنصر وجنده، وقيل هم مشركو قريش حين منعوا النبي ﷺ عن المسجد الحرام.

\isnad{ورجح الطبري:} «وأولى التأويلات... قول من قال: عنى الله عز وجل بقوله... النصارى. وذلك أنهم هم الذين سعوا في خراب بيت المقدس... والدليل على صحة ما قلنا... قيام الحجة بأنه لا قول في معنى هذه الآية إلا أحد الأقوال الثلاثة... وكان معلوماً أن مشركي قريش لم يسعوا قط في تخريب المسجد الحرام... وأخرى أن الآية التي قبل قوله \textit{ومن أظلم} مضت بالخبر عن اليهود والنصارى... فإلحاق الآية بما هو نظير قصة الآية قبلها أولى».

\begin{fawaid}[تخصيص اللفظ العام بالسياق والقرائن]
رغم بقاء بعض الألفاظ على عمومها، يميل الطبري إلى تخصيص اللفظ العام ببعض أفراده إذا وُجدت قرينة قوية. فرجّح أن (الذين منعوا مساجد الله وسعوا في خرابها) هم النصارى (مع بختنصر) وسعيهم في خراب بيت المقدس؛ مستدلاً بالسياق السابق واللاحق الذي يذم اليهود والنصارى، وبأن قريشاً لم تكن تسعى في "خراب" المسجد الحرام بل في عمارته وتفتخر بذلك، وإن كانت قد منعت النبي ﷺ منه. وهذا من دقة توظيف السياق والواقع في الترجيح، مع بقاء العبرة بعموم الآية في أن "كل مانع مصلياً في مسجد وكل ساع في خرابه فهو من المعتدين".
\end{fawaid}

\separator

\section{تأويل (أولئك ما كان لهم أن يدخلوها إلا خائفين...)}
\isnad{قال أبو جعفر:} «وهذا خبر من الله جل ثناؤه عمن منع مساجد الله أن يذكر فيها اسمه أنه قد حرّم عليهم دخول المساجد التي سعوا في تخريبها... إلا على خوف ووجل من العقوبة».

\isnad{وأما قوله:} \begin{ayah}لَهُمْ فِي الدُّنْيَا خِزْيٌ\end{ayah} فإنه يعني بالخزي: الشنار والعار، كالقتل والسبي وأداء الجزية عن صغار. \begin{ayah}وَلَهُمْ فِي الْآخِرَةِ عَذَابٌ عَظِيمٌ\end{ayah}: عذاب جهنم الذي لا يخفف.

\separator

\section{تأويل (ولله المشرق والمغرب فأينما تولوا فثم وجه الله)}
\isnad{قال أبو جعفر:} «لله ملكهما وتدبيرهما... يعني أن له ملك ما بين قطري المشرق والمغرب».

واختلفوا في سبب نزولها ومرادها: فقيل في توجيه القبلة من بيت المقدس للرد على اليهود، وقيل نزلت رخصة للمسافر في صلاة التطوع على راحلته، وقيل نزلت في قوم عميت عليهم القبلة فصلوا، وقيل في صلاة النبي ﷺ على النجاشي.

\isnad{ورجح الطبري أنها عامة:} «فإن الصواب فيه من القول أن يقال إنها جاءت مجيء العموم... محتمل \textit{فأينما تولوا} في حال سيركم في أسفاركم في صلاتكم تطوعاً... ومحتمل \textit{فأينما تولوا} من أرض الله كيف تكونون بها فثَمّ قبلة الله... فإذا كان محتملاً ما ذكرنا من الأوجه لم يكن لأحد أن يزعم أنها ناسخة ولا منسوخة إلا بحجة يجب التسليم لها...».

وهنا وضع قاعدة في النسخ: «لا ناسخ من آي القرآن وأخبار رسول الله ﷺ إلا ما نفى حكماً ثابتاً قد لزم العباد فرضه... ولم يصح واحد من هذين المعنيين بقوله \textit{فأينما تولوا فثم وجه الله} بحجة يجب التسليم لها».

\separator

\section{تأويل (وقالوا اتخذ الله ولدا سبحانه بل له ما في السماوات والأرض)}
\isnad{قال أبو جعفر:} «وهم النصارى الذين زعموا أن عيسى ابن الله. فقال جل ثناؤه مكذباً قيلهم... ومنتفياً مما نحلوه وأضافوه إليه... \textit{سبحانه}: تنزيهاً لله وتبرئاً من أن يكون له ولد... ثم أخبر جل ثناؤه أن له ما في السماوات والأرض ملكاً وخلقاً... ولو كان المسيح ابناً كما زعمتم لم يكن كسائر ما في السماوات والأرض من خلقه وعبيده في ظهور آثار الصنعة فيه».

\separator

\section{تأويل (كل له قانتون)}
اختلفوا في القنوت هنا: فقيل مطيعون، وقيل مقرون بالعبودية، وقيل قيام يوم القيامة.
\isnad{قال أبو جعفر:} «وللقنوت في كلام العرب معان: الطاعة، والقيام، والكف عن الكلام. وأولى معاني القنوت في قوله \textit{كل له قانتون} الطاعة والإقرار لله بالعبودية وبشهادة أجسامهم بما فيها من آثار الصنعة والدلالة على وحدانيته... فإنهم والسماوات والأرض فيهما إما باللسان وإما بالدلالة».

\separator

\section{تأويل (بديع السماوات والأرض)}
\isnad{قال أبو جعفر:} «بديع السماوات والأرض مبدعها، وإنما هو مُفْعِل فصُرف إلى فَعيل... ومعنى المبدع: المنشئ والمُحدث ما لم يسبقه إلى إنشاء مثله أحد... وهو خالقها وموجدها من غير أصل ولا مثال احتذاها عليه. وهذا إعلام من الله عباده أن مما يشهد له بذلك المسيح... وأن الذي ابتدع السماوات والأرض... هو الذي ابتدع المسيح من غير والد بقدرته».

\separator

\section{تأويل (وإذا قضى أمرا فإنما يقول له كن فيكون)}
\isnad{قال أبو جعفر:} «وإذا أحكم أمراً وحتمه، وأصل كل قضاء: الإحكام والفراغ منه... \textit{فإنما يقول له كن فيكون} فيكون ذلك الأمر على ما أمره الله أن يكون وأراده».

وطرح \rawi{الطبري} إشكالاً فلسفياً أثاره المتكلمون: كيف يأمر الله المعدوم بقوله "كن"؟
فذكر أن بعضهم خصّ الآية بالموجودين (مثل أمر بني إسرائيل كونوا قردة خاسئين)، وبعضهم أوّلها بأنها عبارة عن التكوين بلا قول (مثل: قال الجدار فمال)، وردّ \rawi{الطبري} هذا التحريف للكلم عن مواضعه.

\isnad{ورجح الطبري:} «وأولى الأقوال بالصواب... أن يقال هو عام في كل ما قضاه الله ودبره، لأن ظاهر ذلك ظاهر عموم، وغير جائز إحالة الظاهر إلى الباطن من التأويل بغير برهان... فغير جائز أن يكون الشيء مأموراً بالوجود مراداً كذلك إلا وهو موجود (أي وجوده التكويني عقب الأمر مباشرة)».

\begin{fawaid}[الرد على المتكلمين في (كن فيكون) واختراع الإشكالات]
استشكل بعض المتكلمين قوله: \textit{فَإِنَّمَا يَقُولُ لَهُ كُن فَيَكُونُ}، متسائلين: كيف يأمر الله المعدوم قبل وجوده وهو معدوم؟ فحرفوا الآية وزعموا أنها تمثيل لحال التكوين وأن الله لا يقول "كن" حقيقةً. فأجاب الطبري متمسكاً بظاهر القرآن أن الله يأمر الشيء حقيقة بـ "كن"، وأن أمره وإرادته للشيء لا تتقدم وجوده ولا تتأخر عنه، فبمجرد أمره يوجد. والمشكلة أصلاً في افتراض المتكلمين قواعد عقلية بشرية لمحاكمة النصوص، بينما السلف لم يستشكلوا ذلك لعلمهم بكمال قدرة الله المطلقة التي توجد الأشياء بأمره.
\end{fawaid}

\separator

\section{خاتمة اللقاء}
بهذا القدر نكتفي اليوم، ونسأل الله أن ينفعنا بما قرأنا وبما تعلمنا، ونكمل غداً أو في المجلس القادم بإذن الله باقي آيات سورة البقرة. والسلام عليكم ورحمة الله وبركاته.